\documentclass[conference]{IEEEtran}

% ---------- Packages ----------
\usepackage[T1]{fontenc}
\usepackage{cite}
\usepackage{amsmath,amssymb,amsfonts}
\usepackage{graphicx}
\usepackage{booktabs}
\usepackage{multirow}
\usepackage{array}
\usepackage{url}
\usepackage{xcolor}
\usepackage{hyperref}
\usepackage{balance}

% ---------- Hyperref ----------
\hypersetup{
    colorlinks=true,
    linkcolor=blue,
    citecolor=blue,
    urlcolor=blue
}

% ---------- Title and Authors ----------
\title{Your Paper Title: A Concise and Informative Description}

\author{
\IEEEauthorblockN{WEN WEI CHENG\IEEEauthorrefmark{1}}
\IEEEauthorblockA{\IEEEauthorrefmark{1}Department of Computer Science, National Chung Hsing University\\
Taiwan\\
wilber830915@gmail.com}
}

\begin{document}

\maketitle

% ---------- Abstract ----------
\begin{abstract}
Write a concise summary of the research problem, motivation, proposed
method, experimental design, key results, and main contribution.
Avoid citations and unexplained abbreviations where possible.
\end{abstract}

\begin{IEEEkeywords}
Keyword 1, Keyword 2, Keyword 3, Keyword 4
\end{IEEEkeywords}

% ---------- 1. Introduction ----------
\section{Introduction}
Introduce the research background and practical motivation. Clearly define
the problem, explain why it matters, and summarize the limitations of existing
approaches. End this section by listing the main contributions.

The main contributions of this paper are:
\begin{itemize}
    \item Contribution 1: describe the main research contribution.
    \item Contribution 2: describe the proposed method or system.
    \item Contribution 3: describe the evaluation and findings.
\end{itemize}

% ---------- 2. Related Work ----------
\section{Related Work}
Review relevant studies and group them by research theme. Explain how the
proposed work differs from prior approaches. Cite sources using BibTeX,
for example \cite{example2023}.

% ---------- 3. Proposed Method ----------
\section{Proposed Method}
Describe the system architecture, methodology, algorithms, and
implementation details so that the work can be understood and reproduced.

\subsection{System Overview}
Explain the overall workflow and its major components.

% Example figure: upload your image to the Overleaf project and replace
% figures/system_overview.png with its actual filename.
% \begin{figure}[t]
%     \centering
%     \includegraphics[width=\linewidth]{figures/system_overview.png}
%     \caption{Overall architecture of the proposed system.}
%     \label{fig:architecture}
% \end{figure}

\subsection{Mathematical Formulation}
Define the variables before or immediately after using them. For example,
the mean squared error (MSE) can be expressed as
\begin{equation}
    \mathrm{MSE} = \frac{1}{N}\sum_{i=1}^{N}(y_i-\hat{y}_i)^2,
    \label{eq:mse}
\end{equation}
where $N$ is the number of samples, $y_i$ is the ground-truth value, and
$\hat{y}_i$ is the predicted value.

% ---------- 4. Experimental Setup ----------
\section{Experimental Setup}
Describe the dataset, preprocessing, hardware/software environment,
baseline methods, parameter settings, and evaluation metrics.

\subsection{Dataset}
Describe the data source, sample size, features, labels, and data split.

\subsection{Implementation Details}
Report the software versions, hardware specifications, and important
hyperparameters needed for reproducibility.

% ---------- 5. Results and Discussion ----------
\section{Results and Discussion}
Present quantitative and qualitative results. Interpret the findings rather
than only repeating the values in tables.

\begin{table}[t]
\caption{Example quantitative comparison.}
\label{tab:results}
\centering
\begin{tabular}{lcc}
\toprule
Method & Metric 1 & Metric 2 \\
\midrule
Baseline A & 0.00 & 0.00 \\
Baseline B & 0.00 & 0.00 \\
Proposed Method & 0.00 & 0.00 \\
\bottomrule
\end{tabular}
\end{table}

Discuss the strengths, limitations, possible sources of error, and
implications of the results. Refer to tables and figures in the text,
e.g., Table~\ref{tab:results}.

% ---------- 6. Conclusion ----------
\section{Conclusion}
Summarize the problem, proposed approach, and major findings. State the
limitations and identify concrete directions for future work.

% ---------- Acknowledgment (optional) ----------
\section*{Acknowledgment}
The authors would like to thank ...

% ---------- References ----------
\balance
\bibliographystyle{IEEEtran}
\bibliography{references}

\end{document}
